\documentclass[10pt,a4paper]{amsart}
\usepackage{amsmath,amssymb,amsthm,graphicx}
\usepackage{hyperref}
\numberwithin{equation}{section}

\newtheorem{Theorem}{Theorem}[section]
\newtheorem*{Theorem*}{Theorem}
\newtheorem{Corollary}[Theorem]{Corollary}
\newtheorem{Lemma}[Theorem]{Lemma}
\newtheorem{Proposition}[Theorem]{Proposition}
 { \theoremstyle{definition}
\newtheorem{Definition}[Theorem]{Definition}
\newtheorem{Note}[Theorem]{Note}
\newtheorem{Example}[Theorem]{Example}
\newtheorem{Remark}[Theorem]{Remark} }

\def\Var{\mathrm{Var}}
\def\Nov{\mathrm{Nov}}
\def\Com{\mathrm{Com}}
\def\As{\mathrm{As}}
\def\Lie{\mathrm{Lie}}
\def\Perm{\mathrm{Perm}}
\def\Zinb{\mathrm{Zinb}}
\def\Pois{\mathrm{Pois}}
\def\Jord{\mathrm{Jord}}
\def\SJord{\mathrm{SJord}}
\def\pre{\mathrm{pre}}
\def\RB{\mathrm{RB}}
\def\wt {\mathop {\fam 0 wt}\nolimits}
\def\id {\mathop {\fam 0 id}\nolimits}
\def\Der {\mathop {\fam 0 Der}\nolimits}
\def\End {\mathop {\fam 0 End}\nolimits}
\def\oo#1{\mathbin{{}_{(#1)}}}
\def\oon#1{\mathbin{{\circ}_{(#1)}}}

\def\Red#1{\textcolor{red}{#1}}
\def\Blue#1{\textcolor{blue}{#1}}
\title{Every non-associative algebra embeds into an associative algebra through a derivation}
\author{Kolesnikov, Pavel, Mashurov, Farukh, Sartayev, Bauyrzhan}
\date{doi:10.3842/SIGMA.2024.017; CC BY 4.0}
\begin{document}
\maketitle
\setcounter{section}{3}
\setcounter{Theorem}{6}
\noindent\textit{Editorial note.} All delivered mathematical text is the author's own, copied from the pinned SIGMA article source; no statement or proof step was rewritten, reordered or removed. The counted statement is the labeled proposition environment at source lines 812-819 and the counted proof is the proof environment at source lines 826-861.\par\medskip
\noindent\textit{Editorial note.} The proof cites published references of the author's bibliography by key; the cited entries are copied line for line from the author's own bibliography into a retained bibliography fragment, and hyperref is added because the publisher class, not the author preamble, supplied the link command.\par\medskip
\noindent\textit{Editorial note.} The retained preamble of this unit was generated mechanically for this delivery: it copies the author's own package, numbering, macro and theorem declaration block from the pinned source (source lines 9-43), with the publisher class and page setup lines removed.\par\medskip
\noindent\textit{Allowed standard prerequisites (editorial).}\begin{itemize}
\item The free associative algebra generated by an enriched set with an induced well-order. standard construction used at proof lines 827-832. Supplies the ambient algebra in which the embedding is constructed, together with the order used to analyse the products.
\item The rewriting of products of ordered monomials through the derivation induced on the free algebra. derived inside the proof at proof lines 832-859. Shows that the symmetrized products of the images of the elements of the given algebra reproduce its multiplication.
\end{itemize}\medskip
\begin{Proposition}\label{prop:DerAComm}
For every non-associative algebra $V$ with
a multiplication $\nu \colon V\otimes V \to V$
there exists an associative algebra $(A,\cdot )$
with a derivation $d$ such that
$V\subseteq A$ and $\nu (u,v) = d(u)\cdot v + v\cdot d(u)$
for all $u,v\in V$.
\end{Proposition}

\begin{proof}
Let us choose a linear basis $B$ of $V$
equipped with an arbitrary total order $\le $
such that $(B,\le )$ is a well-ordered set.
Then define $F$ to be the free associative algebra
generated by $B^{(\omega )}$.
Induce the order $\le $ on $B^{(\omega )}$
by the following rule:
\[
a^{(n)} \le b^{(m)} \iff (n,a)\le (m,b)\ \text{lexicographically},
\]
and expand it to the words in $B^{(\omega )}$ by the deg-lex rule
(first by length, then lexicographically).

Consider the set of defining relations
\begin{gather*}
S = \bigg\{a^{(n)}b + \sum\limits_{s\ge 1} \binom{n-1}{s}\!
\bigl(a^{(n-s)}b^{(s)} +b^{(s)}a^{(n-s)} \bigr)
+ba^{(n)} - \nu(a,b)^{(n-1)} \mid a,b\in B, n\ge 1\bigg \}.
\end{gather*}
All relations in the set $S$ are obtained
from $a'b+ba'-\nu(a,b)$ by formal derivation
$d\colon x^{(s)}\to x^{(s+1)}$, $x^{(s)}\in B^{(\omega )}$.
Hence, $A=F/(S)$ is a differential associative algebra,
and the map $\varphi \colon V\to A^{(+)}$, $\varphi( v)= v+(S)$, $v\in V$,
preserves the operation, i.e.,
$\varphi(\nu(u,v)) =
d(\varphi(u))\cdot \varphi (v) +\varphi(v)\cdot d(\varphi(u))$ for all $u,v\in V$.

The principal parts of $f\in S$ relative to the order
$\le $ are $a^{(n)}b$, $a,b\in B$, $n\ge 1$.
These words have no compositions of inclusion or intersection, hence, $S$ is a Gr\"obner--Shirshov
basis in $F$ and the images of all variables from $B$ are linearly independent in $A$ since they are $S$-reduced (see, e.g.,~\cite{BC-bull2014} for the
definitions).
Therefore, $\varphi \colon V\to A^{(+)}$ is the desired embedding.
\end{proof}

\begin{thebibliography}{99}
\bibitem{BC-bull2014}
Bokut L.A., Chen Y., Gr\"obner--{S}hirshov bases and their calculation,
 \href{https://doi.org/10.1007/s13373-014-0054-6}{\textit{Bull. Math. Sci.}} \textbf{4} (2014), 325--395, \href{https://arxiv.org/abs/1303.5366}{arXiv:1303.5366}.
\end{thebibliography}
\end{document}
